% !TeX root = ../main.tex
\SourcePage{24}
$\tilde{\mathfrak{g}}\otimes K_x$ 的：

\begin{equation}\tag{19}
\mathbf{1}\mapsto\mathbf{1},\qquad
u\mapsto gug^{-1}+\operatorname{Res}_x c(u,g^{-1}dg)\cdot\mathbf{1},
\quad u\in\mathfrak{g}\otimes K_x .
\end{equation}

特别地，我们有 Harish-Chandra 对 $\bigl(\tilde{\mathfrak{g}}\otimes K_x, G(\mathcal{O}_x)\bigr)$。
它是 $\bigl(\mathfrak{g}\otimes K_x, G(\mathcal{O}_x)\bigr)$ 被 $\mathbb{C}$ 的中心扩张。记
\[
\mathfrak{z}_x=\mathfrak{z}_{\mathfrak{g}}(\mathcal{O}_x):=D'_{(\mathfrak{g}\otimes K_x,G(\mathcal{O}_x))},
\]
其中 $D'$ 的含义同 1.2.5。

\bdp{2.5.2}\emph{定理.}（[FF92]）。
\begin{enumerate}[label=(\roman*),leftmargin=2.2em]
\item 代数 $\mathfrak{z}_x$ 交换。
\item $\bigl(\tilde{\mathfrak{g}}\otimes K_x, G(\mathcal{O}_x)\bigr)$ 满足扭局部量子化条件（见 1.2.5）。
  也就是说，典范态射 $\sigma_{\mathfrak{z}_x}: \operatorname{gr}\mathfrak{z}_x\to\mathfrak{z}^{cl}_x$ 是同构。\hfill$\square$
\end{enumerate}

\emph{注.} 定理的论断 (i) 在 [FF92] 中对任意代数群 $G$ 有证明，
也对由 $\mathfrak{g}$ 上对称不变双线性型定义的 $\mathfrak{g}\otimes K_x$ 的任意中心扩张有证明。
此外，[FF92] 还证明：完备化扭泛包络代数 $\overline{U}'(\mathfrak{g}\otimes K_x)$ 的中心的
$\pi_0(G(K_x))$-不变部分满射到 $\mathfrak{z}_x$ 上。
Feigin–Frenkel 关于 (i) 的证明的一个版本将在 2.9.3–2.9.5 中给出。
我们已在 2.4.2 末尾的注中解释了其证明的“经典极限”。

\bdp{2.5.3} 线丛 $\omega_{\operatorname{Bun}_G}$ 在 $\operatorname{Bun}_{G,\underline{x}}$ 上定义一个 $G(\mathcal{O}_x)$-等变丛。
$\operatorname{Bun}_{G,\underline{x}}$ 上的 $(\mathfrak{g}\otimes K_x, G(\mathcal{O}_x))$-作用典范地提升为
该线丛上的 $(\tilde{\mathfrak{g}}\otimes K_x, G(\mathcal{O}_x))$-作用，使得 $\mathbf{1}$ 的作用为乘以 2。
事实上，由 2.1.1 有 $\omega_{\operatorname{Bun}_G}=f^*(\det R\Gamma)$，
其中 $f: \operatorname{Bun}_G\to\operatorname{Bun}_{SL(\mathfrak{g})}$ 由伴随表示 $G\to SL(\mathfrak{g})$ 诱导，
而 $\det R\Gamma$ 是 $\operatorname{Bun}_{SL(\mathfrak{g})}$ 上的行列式线丛。
另一方面，众所周知（例如见 [BLa94]），$\det R\Gamma$ 到 $\operatorname{Bun}_{SL_n,\underline{x}}$ 的拉回上
配有 $sl_n(K_x)$ 的 Kac–Moody 扩张（水平 $-1$）的作用。

\SourcePage{25}
\emph{注.} 事实上，该扩张的作用可积分为 $SL_n(K_x)$ 的某个中心扩张的作用（例如见 [BLa94]）。
因此得到典范的中心扩张
\begin{equation}\tag{20}
0\to\mathbb{G}_m\to\hat{G}(K_x)\to G(K_x)\to 0
\end{equation}
它在 $\omega_{\operatorname{Bun}_G}$ 到 $\operatorname{Bun}_{G,\underline{x}}$ 的拉回上的作用
使得 $\lambda\in\mathbb{G}_m$ 的作用为乘以 $\lambda$。
扩张 $0\to\mathbb{C}\to\tilde{\mathfrak{g}}\otimes K_x\to\mathfrak{g}\otimes K_x\to 0$
是对应于 (20) 的李代数扩张的一半。
在第 4 章中，我们将引入 $\omega_{\operatorname{Bun}_G}$ 的一个平方根\footnote{该平方根与扩张 (21) 依赖于 $\omega_X$ 的平方根的选取。}（Pfaffian 丛）以及中心扩张
\begin{equation}\tag{21}
0\to\mathbb{G}_m\to\tilde{G}(K_x)\to G(K_x)\to 0
\end{equation}
（见 4.4.8），它是 (20) 的一个平方根。
这些平方根对我们比 $\omega_{\operatorname{Bun}_G}$ 与 (20) 更重要，因此我们不给出 $\tilde{G}(K_x)$ 的精确定义。

\bdp{2.5.4} 由 2.5.3 与 1.2.5，我们有滤代数的一个典范态射
\[
h_x: \mathfrak{z}_x\to\Gamma(\operatorname{Bun}_G,D').
\]
在 2.7.5 中我们将证明下述定理。

\bdp{2.5.5}\emph{定理.} 我们的数据满足扭整体量子化条件（见 1.2.5）。\hfill$\square$

如 1.2.3 所解释，由于局部与整体量子化条件都满足，
我们得到 2.2.5 的量子化问题的一个解 $\mathfrak{z}^{(x)}(X)$：
置 $\mathfrak{z}^{(x)}(X)=h_x(\mathfrak{z}_x)$，并给 $\mathfrak{z}^{(x)}(X)$ 配上
由 $\Gamma(\operatorname{Bun}_G,D')$ 上的滤链诱导的滤链
（2.5.5 意味着它也可由 $\mathfrak{z}_x$ 上的滤链诱导）；
于是象征映射把 $\operatorname{gr}\mathfrak{z}^{(x)}(X)$ 与 $h_x^{cl}(\mathfrak{z}^{cl}_x)$ 等同，
且由 2.4.3 有
$h_x^{cl}(\mathfrak{z}^{cl}_x)=h_X^{cl}(\mathfrak{z}^{cl}_x(X))\simeq\mathfrak{z}^{cl}(X)$。

\SourcePage{26}
定理 2.5.5 的证明基于 2.2.5 中量子化问题的解的第二种构造；
它还表明 $\mathfrak{z}^{(x)}(X)$ 不依赖于 $x$。

\emph{注.} 若 $G$ 单连通，则 2.5.5 立即由 2.2.4(v) 得到。

\bdsec{2.6. \texorpdfstring{$\mathcal{D}_X$}{D\_X}-概形的一般性讨论}

\bdp{2.6.1} 设 $X$ 为任意光滑连通代数簇。
一个 $\mathcal{D}_X$-概形是指带有沿 $X$ 的平坦联络的 $X$-概形。
在 $X$ 上仿射的 $\mathcal{D}_X$-概形是交换 $\mathcal{D}_X$-代数的谱
（交换 $\mathcal{D}_X$-代数即带有平坦联络的拟凝聚 $\mathcal{O}_X$-代数）。
$\mathcal{O}_X$-代数 $\mathcal{A}$ 在 $x\in X$ 处的纤维记作 $\mathcal{A}_x$；
这尤其适用于 $\mathcal{D}_X$-代数。
对 $\mathbb{C}$-代数 $C$，记 $C_X$ 为相应的“常值” $\mathcal{D}_X$-代数
（即 $C_X$ 是 $C\otimes\mathcal{O}_X$ 配以显然的联络）。

\bdp{2.6.2}\emph{命题.} 设 $X$ 完备。
\begin{enumerate}[label=(\roman*),leftmargin=2.2em]
\item 函子 $C\rightsquigarrow C_X$ 有左伴随函子：对 $\mathcal{D}_X$-代数 $\mathcal{A}$，
  存在 $\mathbb{C}$-代数 $H_\nabla(X,\mathcal{A})$，使得
  \begin{equation}\tag{22}
  \operatorname{Hom}(\mathcal{A},C_X)=\operatorname{Hom}(H_\nabla(X,\mathcal{A}),C)
  \end{equation}
  对任意 $\mathbb{C}$-代数 $C$ 成立。
\item 典范投影 $\theta_{\mathcal{A}}: \mathcal{A}\to H_\nabla(X,\mathcal{A})_X$ 满。
  于是 $H_\nabla(X,\mathcal{A})_X$ 是 $\mathcal{A}$ 的最大“常值”商 $\mathcal{D}_X$-代数。
  特别地，对任意 $x\in X$，态射
  $\theta_{\mathcal{A}_x}: \mathcal{A}_x\to(H_\nabla(X,\mathcal{A})_X)_x=H_\nabla(X,\mathcal{A})$ 满。
\end{enumerate}

\emph{注.} (i) 这里的代数不假定交换、结合等。我们将在 $\mathcal{A}$ 交换的情形下需要该命题。

(ii) 设 $\mathcal{A}$ 交换（这是“交换结合含幺”的缩写）。
则按命题的论断 (ii)，$H_\nabla(X,\mathcal{A})$ 交换。
若 $C$ 也假定交换，则 (22) 只是说
$\operatorname{Spec} H_\nabla(X,\mathcal{A})$ 是 $\operatorname{Spec}\mathcal{A}$ 的水平截面概形。

\SourcePage{27}
从几何观点看，这样的概形存在且仿射是清楚的：
$\operatorname{Spec}\mathcal{A}$ 的所有截面构成仿射概形 $S$
（这里用到 $X$ 的完备性；否则 $S$ 将是归纳概形，见下一个注），
而水平截面构成 $S$ 的闭子概形。
$\theta_{\mathcal{A}_x}$ 与 $\theta_{\mathcal{A}}$ 的满性意味着态射
$\operatorname{Spec} H_\nabla(X,\mathcal{A})\to\operatorname{Spec}\mathcal{A}_x$
与 $X\times\operatorname{Spec} H_\nabla(X,\mathcal{A})\to\operatorname{Spec}\mathcal{A}$ 是闭浸入。

(iii) 若 $X$ 任意（不必完备），则由 (22) 定义的 $H_\nabla(X,\mathcal{A})$
可表为代数关于某个有向满射族的投射极限。
于是若 $\mathcal{A}$ 交换，则 $\operatorname{Spec}\mathcal{A}$ 的水平截面空间是归纳仿射归纳概形\footnote{从几何观点看这也很清楚。事实上，水平截面构成所有截面空间 $S_X$ 中的闭子空间。若 $X$ 仿射，$S_X$ 当然是归纳概形。一般情形下 $X$ 可被开仿射子概形 $U_1,\dots,U_n$ 覆盖；于是 $S_X$ 是诸 $S_{U_i}$ 之积的闭子空间。}。

\emph{证明.} (a) 记 $\mathcal{M}(X)$ 为 $\mathcal{D}_X$-模的范畴，
$\mathcal{M}_{const}(X)$ 为常值 $\mathcal{D}_X$-模的满子范畴，
即同构于 $V\otimes\mathcal{O}_X$（$V$ 为某向量空间）的 $\mathcal{D}_X$-模
（实际上函子 $V\mapsto V\otimes\mathcal{O}_X$ 是向量空间范畴与 $\mathcal{M}_{const}(X)$ 之间的等价）。
我们断言嵌入 $\mathcal{M}_{const}(X)\to\mathcal{M}(X)$ 有左伴随函子，
即对 $\mathcal{F}\in\mathcal{M}(X)$ 存在 $\mathcal{F}_\nabla\in\mathcal{M}_{const}(X)$，使得
$\operatorname{Hom}(\mathcal{F},\mathcal{E})=\operatorname{Hom}(\mathcal{F}_\nabla,\mathcal{E})$
对 $\mathcal{E}\in\mathcal{M}_{const}(X)$ 成立。
只需对凝聚的 $\mathcal{F}$ 构造 $\mathcal{F}_\nabla$。
此时 $\mathcal{F}_\nabla:=(\operatorname{Hom}_{\mathcal{D}_X}(\mathcal{F},\mathcal{O}_X))^*\otimes\mathcal{O}_X$
（这里用到 $\dim\operatorname{Hom}_{\mathcal{D}_X}(\mathcal{F},\mathcal{O}_X)<\infty$，因为 $X$ 完备）。

(b) 由于 $\mathcal{O}_X$ 是不可约 $\mathcal{D}_X$-模，常值 $\mathcal{D}_X$-模的 $\mathcal{D}_X$-子模是常值的。
于是自然态射 $\mathcal{F}\to\mathcal{F}_\nabla$ 满。

(c) 若 $\mathcal{A}$ 是 $\mathcal{D}_X$-代数，$\mathcal{I}$ 是 $\mathcal{A}$ 中由
$\operatorname{Ker}(\mathcal{A}\to\mathcal{A}_\nabla)$ 生成的理想，
则 $\mathcal{A}/\mathcal{I}$ 是常值 $\mathcal{D}_X$-模 $\mathcal{A}_\nabla$ 的商。
所以 $\mathcal{A}/\mathcal{I}$ 常值，即 $\mathcal{A}/\mathcal{I}=H_\nabla(X,\mathcal{A})\otimes\mathcal{O}_X$，
其中 $H_\nabla(X,\mathcal{A})$ 为某向量空间。
$\mathcal{A}/\mathcal{I}$ 是 $\mathcal{D}_X$-代数，故 $H_\nabla(X,\mathcal{A})$ 是代数。
显然它满足 (22)。\hfill$\square$

\SourcePage{28}
\emph{注.} 偏几何的读者可把上面的注 (ii) 看作交换代数情形下命题的证明。
然而在 2.7.4 中，我们将在 $\mathcal{A}$ 交换而 $C=\Gamma(\operatorname{Bun}_G,D')$ 并不显然交换的情形下应用 (22)。
此时只需注意：态射 $\mathcal{A}\to C\otimes\mathcal{O}_X$ 的像形如 $C'\otimes\mathcal{O}_X$
（见命题证明的 (b) 部分），且 $C'$ 交换，因为 $\mathcal{A}$ 交换。
也可把 (22) 用于 $C:=\Gamma(\operatorname{Bun}_G,D')$ 中由所有 $x\in X$ 的态射
$h_x: \mathfrak{z}_x\to\Gamma(\operatorname{Bun}_G,D')$ 的像生成的子代数
（这个 $C$ “显然”交换；见 2.9.1）。
事实上，用 2.2.4(v) 可以证明 $\Gamma(\operatorname{Bun}_G,D')$ 交换
（由 2.2.4(v) 可推出：对任意连通分支
$\operatorname{Bun}_G^\gamma\subset\operatorname{Bun}_G$ 与任意 $x\in X$，
由 $h_x$ 诱导的态射 $\mathfrak{z}_x\to\Gamma(\operatorname{Bun}_G^\gamma,D')$ 满）。

\bdp{2.6.3} 本小节中所有代数均假定交换。
遗忘函子 $\{\mathcal{D}_X\text{-代数}\}\to\{\mathcal{O}_X\text{-代数}\}$ 有显然的左伴随函子 $\mathcal{J}$
（$\mathcal{J}\mathcal{A}$ 是由 $\mathcal{O}_X$-代数 $\mathcal{A}$ 生成的 $\mathcal{D}_X$-代数）。
我们断言 $\operatorname{Spec}\mathcal{J}\mathcal{A}$ 正是
\emph{$\operatorname{Spec}\mathcal{A}$ 的截面的 $\infty$-射流概形}。
特别地，这意味着 $\operatorname{Spec}(\mathcal{J}\mathcal{A})_x$ 的 $\mathbb{C}$-点
与截面 $\operatorname{Spec}\mathcal{O}_x\to\operatorname{Spec}\mathcal{A}$ 之间有典范的一一对应
（其中 $\mathcal{O}_x$ 是 $x$ 处局部环的形式完备化）。
更确切地说，我们要构造函子性的双射
\begin{equation}\tag{23}
\operatorname{Hom}_{\mathcal{O}_X}(\mathcal{J}\mathcal{A},\mathcal{B})
\overset{\sim}{\longrightarrow}
\operatorname{Hom}_{\mathcal{O}_X}(\mathcal{A},\hat{\mathcal{B}})
\end{equation}
其中 $\mathcal{B}$ 是（拟凝聚）$\mathcal{O}_X$-代数，$\hat{\mathcal{B}}$ 是
$\mathcal{O}_X\otimes_{\mathbb{C}}\mathcal{B}$ 关于理想
$\operatorname{Ker}(\mathcal{O}_X\otimes_{\mathbb{C}}\mathcal{B}\to\mathcal{B})$ 的完备化。
这里 $\hat{\mathcal{B}}$ 配有由态射
$\mathcal{O}_X\to\mathcal{O}_X\otimes_{\mathbb{C}}\mathcal{B}$（$a\mapsto a\otimes 1$）给出的
$\mathcal{O}_X$-代数结构。
让我们暂时去掉 $\mathcal{D}_X$-代数定义中的拟凝聚假定。
则 $\hat{\mathcal{B}}$ 是 $\mathcal{D}_X$-代数
（$\hat{\mathcal{B}}$ 上的联络来自 $\mathcal{O}_X\otimes_{\mathbb{C}}\mathcal{B}$ 上的联络，
该联络使 $1\otimes\mathcal{B}$ 的截面水平）。
于是 $\operatorname{Hom}_{\mathcal{O}_X}(\mathcal{A},\hat{\mathcal{B}})
=\operatorname{Hom}_{\mathcal{D}_X}(\mathcal{J}\mathcal{A},\hat{\mathcal{B}})$，
而要构造 (23)，只需构造函子性的双射
\SourcePage{29}
\begin{equation}\tag{24}
\operatorname{Hom}_{\mathcal{O}_X}(\mathcal{R},\mathcal{B})
\longleftarrow
\operatorname{Hom}_{\mathcal{D}_X}(\mathcal{R},\hat{\mathcal{B}})
\end{equation}
对任意 $\mathcal{D}_X$-代数 $\mathcal{R}$ 与 $\mathcal{O}_X$-代数 $\mathcal{B}$ 成立
（即证明函子 $\mathcal{B}\mapsto\hat{\mathcal{B}}$ 是遗忘函子
$\{\mathcal{D}_X\text{-代数}\}\to\{\mathcal{O}_X\text{-代数}\}$ 的\emph{右}伴随）。
映射 (24) 来自显然的态射 $\hat{\mathcal{B}}\to\mathcal{B}$。
读者容易证明 (24) 是双射。

对 $\mathcal{D}_X$-代数 $\mathcal{A}$ 与 $\mathbb{C}$-代数 $C$，有
\[
\operatorname{Hom}_{\mathcal{D}_X\text{-alg}}(\mathcal{J}\mathcal{A},C\otimes\mathcal{O}_X)
=\operatorname{Hom}_{\mathcal{O}_X\text{-alg}}(\mathcal{A},C\otimes\mathcal{O}_X).
\]
这意味着：\emph{典范态射 $\operatorname{Spec}\mathcal{J}\mathcal{A}\to\operatorname{Spec}\mathcal{A}$
把 $\operatorname{Spec}\mathcal{J}\mathcal{A}$ 的水平截面归纳概形
与 $\operatorname{Spec}\mathcal{A}$ 的（所有）截面的归纳概形等同}。
若 $X$ 完备，则由 2.6.2，这些空间实际上是概形。

最后我们指出，本小节的结果可按显然的方式整体化。
遗忘函子 $\{\mathcal{D}_X\text{-概形}\}\to\{X\text{-概形}\}$ 有右伴随函子
$\mathcal{J}: \{X\text{-概形}\}\to\{\mathcal{D}_X\text{-概形}\}$。
对 $X$-概形 $Y$，$\mathcal{J}Y$ 是 $Y$ 的截面的 $\infty$-射流概形。
对 $\mathcal{O}_X$-代数 $\mathcal{A}$，有 $\mathcal{J}\operatorname{Spec}\mathcal{A}=\operatorname{Spec}\mathcal{J}\mathcal{A}$。
典范态射 $\mathcal{J}Y\to Y$ 把 $\mathcal{J}Y$ 的水平截面空间\footnote{最一般的情形下，“空间”指“函子 $\{\mathbb{C}\text{-代数}\}\to\{\text{集合}\}$”。}
与 $Y$ 的（所有）截面空间等同。
若 $X$ 完备且 $Y$ 拟射影，则我们的空间是概形。

\bdp{2.6.4} 设 $(\mathfrak{l},P)$ 是 1.2.6 意义下的 Harish-Chandra 对
（故 $P$ 可为任意仿射群概形；我们不假定它有限型\footnote{由下面的定义可见，$\operatorname{Lie} P$ 在 $\mathfrak{l}$ 中的余维有限（等于 $\dim X$）。}）。

\emph{定义.} $X$ 上的 $(\mathfrak{l},P)$-结构是指态射 $\pi: X^{\wedge}\to X$
连同 $(\mathfrak{l},P)$ 在 $X^{\wedge}$ 上的作用，使得
\begin{enumerate}[label=(\roman*),leftmargin=2.2em]
\item $X^{\wedge}$ 是 $X$ 上的 $P$-torsor。
\item $\mathfrak{l}$ 在 $X^{\wedge}$ 上的作用形式自由且传递，即给出同构
  $\mathfrak{l}\hat{\otimes}\mathcal{O}_{X^{\wedge}}\overset{\sim}{\to}\Theta_{X^{\wedge}}$。
\end{enumerate}

\SourcePage{30}
\emph{注.} 设 $L$ 是李代数为 $\mathfrak{l}$、$L_{red}=P$ 的群归纳概形（见 7.11.2(v)）。
% [Tutor Note: 原书 p.30 此处拼作「homogenuous」，疑为「homogeneous」之笔误，底本如此。]
考虑齐性空间 $P\backslash L=\operatorname{Spf}\mathcal{O}$，
其中 $\mathcal{O}=\mathcal{O}_{(\mathfrak{l},P)}=(U\mathfrak{l}/(U\mathfrak{l})\mathfrak{p})^*$。
取 $x\in X$ 并选 $x^{\wedge}\in\pi^{-1}(x)$。
映射 $L\to X^{\wedge}$，$l\mapsto lx^{\wedge}$，给出态射
$\alpha_{x^{\wedge}}: \operatorname{Spf}\mathcal{O}\to X$，
它把 $\operatorname{Spf}\mathcal{O}$ 与 $x$ 的形式邻域等同。
对 $l\in L$，$a\in\operatorname{Spf}\mathcal{O}$，有 $\alpha_{lx^{\wedge}}(a)=\alpha_{x^{\wedge}}(al)$。
注意：若 $P$ 在 $\mathcal{O}$ 上的作用忠实，则 $x^{\wedge}$ 由 $\alpha_{x^{\wedge}}$ 唯一确定。

\bdp{2.6.5}\emph{例.} 置 $\mathcal{O}=\mathcal{O}_n:=\mathbb{C}[[t_1,..,t_n]]$。
$\mathbb{C}$-代数 $\mathcal{O}$ 的自同构群自然地是 $\mathbb{C}$ 上一个仿射群概形
$\operatorname{Aut}^0\mathcal{O}$ 的 $\mathbb{C}$-点群。
记 $\operatorname{Aut}\mathcal{O}$ 为这样的群归纳概形：
对任意 $\mathbb{C}$-代数 $R$，$(\operatorname{Aut}\mathcal{O})(R)$ 是拓扑 $R$-代数
$R\hat{\otimes}\mathcal{O}=R[[t_1,..,t_n]]$ 的自同构群。
于是 $\operatorname{Aut}^0\mathcal{O}$ 是 $\operatorname{Aut}\mathcal{O}$ 的群子概形；
事实上 $\operatorname{Aut}^0\mathcal{O}=(\operatorname{Aut}\mathcal{O})_{red}$。
有 $\operatorname{Lie}\operatorname{Aut}\mathcal{O}=\operatorname{Der}\mathcal{O}$，
$\operatorname{Lie}\operatorname{Aut}^0\mathcal{O}=\operatorname{Der}^0\mathcal{O}
:=\mathfrak{m}_{\mathcal{O}}\cdot\operatorname{Der}\mathcal{O}$。
因此 $\operatorname{Aut}\mathcal{O}$ 是对应于 Harish-Chandra 对
$\operatorname{Aut}^{HC}\mathcal{O}:=(\operatorname{Der}\mathcal{O},\operatorname{Aut}^0\mathcal{O})$ 的群归纳概形。
我们滥用记号，将写 $\operatorname{Aut}\mathcal{O}$ 而不用 $\operatorname{Aut}^{HC}\mathcal{O}$。

正如 Gelfand 与 Kazhdan 所解释（见 [GK]、[GKF] 与 [BR]），
任意 $n$ 维光滑代数簇 $X$ 带有典范的\footnote{事实上，$X$ 上的 $\operatorname{Aut}\mathcal{O}$-结构在同构意义下唯一（这由 2.6.4 的注得出）。}
$\operatorname{Aut}\mathcal{O}$-结构。
空间 $X^{\wedge}=X^{\wedge}_{can}$ 是 $X$ 上“形式坐标系”的空间。
换言之，$X^{\wedge}$ 的一个 $\mathbb{C}$-点是微分非零的态射 $\operatorname{Spec}\mathcal{O}\to X$，
而 $X^{\wedge}$ 的一个 $R$-点是 $R$-态射
$\alpha: \operatorname{Spec}(R\hat{\otimes}\mathcal{O})\to X\otimes R$，
其微分在 $\operatorname{Spec} R$ 的每一点上非零。
群归纳概形 $\operatorname{Aut}\mathcal{O}$ 以显然的方式作用在 $X^{\wedge}$ 上，
且有投影 $\pi: X^{\wedge}\to X$，$\alpha\mapsto\alpha(0)$。
容易看出 $X^{\wedge}$（连同上列结构）是 $X$ 上的 $\operatorname{Aut}\mathcal{O}$-结构。

我们将用到 $n=1$ 时的典范 $\operatorname{Aut}\mathcal{O}_n$-结构，
即 $X$ 为曲线时，此时 $\mathcal{O}=\mathbb{C}[[t]]$。
这里群 $\operatorname{Aut}\mathcal{O}$ 如下。
有满同态
$\operatorname{Aut}^0\mathcal{O}\to\operatorname{Aut}(t\mathcal{O}/t^2\mathcal{O})=\mathbb{G}_m$，
\SourcePage{31}
我们称之为 $\operatorname{Aut}^0\mathcal{O}$ 的\emph{标准特征标}；其核是投射幂幺的。
对 $\mathbb{C}$-代数 $R$，$R[[t]]$ 的自同构由 $t\mapsto\sum_i c_i t^i$ 定义，
其中 $c_1\in R^*$，$c_0$ 幂零。
于是 $\operatorname{Aut}\mathcal{O}$ 是概形
$\operatorname{Spec}\mathbb{C}[c_0,c_1,c_1^{-1},c_2,c_3,\dots]/(c_0^k)$（$k\in\mathbb{N}$）的并。
$\operatorname{Aut}^0\mathcal{O}$ 是由 $c_0=0$ 定义的 $\operatorname{Aut}\mathcal{O}$ 的群子概形。

另一些 $(\mathfrak{l},P)$-结构的例子可在 ?? 中找到。
% [Tutor Note: 原书 p.31 此处引用印作「??」，底本如此，保留原样。]

\bdp{2.6.6} 设 $X$ 是带有 $(\mathfrak{l},P)$-结构 $X^{\wedge}$ 的代数簇
（我们将在下述情形下应用这些构造：$X$ 为曲线，$\mathfrak{l}=\operatorname{Der}\mathcal{O}$，
$P=\operatorname{Aut}^0\mathcal{O}$（或 $\operatorname{Aut}^0\mathcal{O}$ 的某个覆叠），
$\mathcal{O}:=\mathbb{C}[[t]]$）。
记 $\mathcal{M}(X,\mathcal{O})$ 为 $X$ 上 $\mathcal{O}$-模的范畴，
$\mathcal{M}^l(X)$ 为左 $\mathcal{D}$-模的范畴。
对 $\mathcal{F}_X\in\mathcal{M}(X,\mathcal{O})$，其到 $X^{\wedge}$ 的拉回
$\mathcal{F}_{X^{\wedge}}$ 是 $X^{\wedge}$ 上的 $P$-等变 $\mathcal{O}$-模。
若 $\mathcal{F}_X$ 实际上是左 $\mathcal{D}_X$-模，则 $\mathcal{F}_{X^{\wedge}}$ 还是 $\mathfrak{l}$-等变的
（因为由 2.6.4(ii)，$\mathcal{O}_{X^{\wedge}}$-模上的 $\mathfrak{l}$-作用等同于平坦联络）。
函子
$\mathcal{M}(X,\mathcal{O})\to\{X^{\wedge}\text{ 上的 }P\text{-等变 }\mathcal{O}\text{-模}\}$、
$\mathcal{M}^l(X)\to\{X^{\wedge}\text{ 上的 }(\mathfrak{l},P)\text{-等变 }\mathcal{O}\text{-模}\}$
是张量范畴的等价。

有忠实正合张量函子
\begin{equation}\tag{25}
\mathcal{M}(P)\longrightarrow\mathcal{M}(X,\mathcal{O}),\qquad
\mathcal{M}(\mathfrak{l},P)\longrightarrow\mathcal{M}^l(X)
\end{equation}
它们把表示 $V$ 送到 $\mathcal{O}_X$-模或 $\mathcal{D}_X$-模 $V_X$，
使得 $V_{X^{\wedge}}$ 等于 $V\otimes\mathcal{O}_{X^{\wedge}}$
（$P$-模或 $(\mathfrak{l},P)$-模的张量积）。
换言之，$\mathcal{O}_X$-模 $V_X$ 是 $V$ 用 $P$-torsor $X^{\wedge}$ 扭出的结果。
因此任意带 $P$-作用的代数 $A$ 给出 $\mathcal{O}_X$-代数 $A_X$；
若 $A$ 实际上带有 $(\mathfrak{l},P)$-作用，则 $A_X$ 是 $\mathcal{D}_X$-代数。
类似地，任意带 $P$-作用的概形 $H$（简称\emph{$P$-概形}）给出 $X$-概形 $H_X$。
若 $H$ 实际上是 $(\mathfrak{l},P)$-概形，则 $H_X$ 是 $\mathcal{D}_X$-概形。
有 $(\operatorname{Spec} A)_X=\operatorname{Spec}(A_X)$。

\emph{注.} (i) 函子 $\mathcal{M}(\mathfrak{l},P)\longrightarrow\mathcal{M}^l(X)$
与 $X^{\wedge}$ 上 $(\mathfrak{l},P)$-作用的局部化函子 $\Delta$ 一致（见 1.2.4）。

(ii) 函子 (25) 有右伴随，它把 $\mathcal{O}_X$-模或 $\mathcal{D}_X$-模 $\mathcal{F}_X$ 送到向量空间
\SourcePage{32}
$\Gamma(X^{\wedge},\mathcal{F}_{X^{\wedge}})$，其上配有显然的 $P$-模或 $(\mathfrak{l},P)$-模结构。
若考虑代数而非模，同样的伴随性成立。

(iii) 设 $C$ 是 $\mathbb{C}$-代数；把 $C$ 看作带平凡 $\operatorname{Aut}\mathcal{O}$-作用的 $(\mathfrak{l},P)$-代数。
则 $C_X$ 是 2.6.1 中的“常值” $\mathcal{D}_X$-代数。

\bdp{2.6.7} 遗忘函子 $\{(\mathfrak{l},P)\text{-代数}\}\to\{P\text{-代数}\}$
有左伴随（诱导）函子 $\mathcal{J}$。
对 $P$-代数 $A$ 有典范同构
\begin{equation}\tag{26}
(\mathcal{J}A)_X=\mathcal{J}(A_X).
\end{equation}
事实上，自然的 $\mathcal{O}_X$-代数态射 $A_X\to(\mathcal{J}A)_X$ 诱导
$\mathcal{D}_X$-代数态射 $\mathcal{J}(A_X)\to(\mathcal{J}A)_X$。
为证它是同构，用 $\mathcal{J}$ 与 $A\mapsto A_X$ 的伴随性（见 2.6.3 与 2.6.6 的注 (ii)）。

上面陈述的几何版本如下。
遗忘函子 $\{(\mathfrak{l},P)\text{-概形}\}\to\{P\text{-概形}\}$ 有右伴随函子\footnote{对仿射概形，这只是上面关于 $P$-代数的陈述的改写。一般情形不能立即归结为仿射情形（$P$-概形未必有 $P$-不变仿射子概形的覆盖），但仿射情形已够我们使用。} $\mathcal{J}$。
对 $P$-代数 $A$ 有 $\mathcal{J}(\operatorname{Spec} A)=\operatorname{Spec}\mathcal{J}(A)$。
对任意 $P$-概形 $H$ 有 $(\mathcal{J}H)_X=\mathcal{J}(H_X)$。

\bdsec{2.7. 量子化 II}
从现在起 $\mathcal{O}:=\mathbb{C}[[t]]$，$K:=\mathbb{C}((t))$。

\bdp{2.7.1} 先考虑“经典”图景。
概形 $\operatorname{Hitch}_x$（$x\in X$）是在 $X$ 上仿射的 $\mathcal{D}_X$-概形
$\operatorname{Hitch}=\mathcal{J}C_{\omega_X}$ 的纤维；
记 $\mathfrak{z}^{cl}$ 为相应的 $\mathcal{D}_X$-代数。
由 2.6.3，投影 $\operatorname{Hitch}\to C_{\omega_X}$ 把 $\operatorname{Hitch}$ 的水平截面概形
与 $\operatorname{Hitch}(X)$ 等同。换言之
\[
\mathfrak{z}^{cl}(X)=H_\nabla(X,\mathfrak{z}^{cl}),
\]
而 2.4.3 中的投影 $\theta_x^{cl}: \mathfrak{z}_x^{cl}\to\mathfrak{z}^{cl}(X)$
正是命题 2.6.2(ii) 中的典范态射 $\theta_{\mathfrak{z}^{cl}_x}$。

\SourcePage{33}
把 $C$ 经由标准特征标 $\operatorname{Aut}^0\mathcal{O}\to\mathbb{G}_m$
看作 $\operatorname{Aut}^0\mathcal{O}$-概形（见 2.6.5）。
$X$-概形 $C_{\omega_X}$ 与 $C$ 的 $X^{\wedge}$-扭重合。
因此同构 (26) 诱导典范同构
\[
\mathfrak{z}^{cl}=\mathfrak{z}^{cl}_{\mathfrak{g}}(\mathcal{O})_X,
\]
其中 $\mathfrak{z}^{cl}_{\mathfrak{g}}(\mathcal{O})$ 是 $\operatorname{Aut}\mathcal{O}$-代数
$\mathcal{J}(\operatorname{Sym}\mathfrak{g})^G$，
而 $(\operatorname{Sym}\mathfrak{g})^G$ 上的 $\operatorname{Aut}^0\mathcal{O}$-作用
来自与 $(\operatorname{Sym}\mathfrak{g})^G$ 的分次所诱导的作用相反的 $\mathbb{G}_m$-作用（参看 2.2.2 的末尾）。

\bdp{2.7.2} 让我们转向“量子”图景。
置 $\mathfrak{z}_{\mathfrak{g}}(\mathcal{O}):=D'_{(\mathfrak{g}\otimes K,G(\mathcal{O}))}$。
这是交换代数（见 2.5.2(i)）。
$\operatorname{Aut}\mathcal{O}$ 作用在 $\mathfrak{z}_{\mathfrak{g}}(\mathcal{O})$ 上，
因为 $\mathfrak{z}_{\mathfrak{g}}(\mathcal{O})$ 是扭真空模 $\mathit{Vac}'$ 的自同态代数（见 1.2.5），
而 $\operatorname{Aut}\mathcal{O}$ 作用在 $\mathit{Vac}'$ 上。
（后一作用由两条性质刻画：它与 $\operatorname{Aut}\mathcal{O}$ 在 $\tilde{\mathfrak{g}}\otimes K$ 上的自然作用相容，
且真空向量不变；$\operatorname{Aut}\mathcal{O}$ 在 $\tilde{\mathfrak{g}}\otimes K$ 上的作用按拓扑意义理解，
即 $\operatorname{Aut}(\mathcal{O}\hat{\otimes}R)$ 作用在 $\tilde{\mathfrak{g}}\otimes K\hat{\otimes}R$ 上，
对任意交换 $\mathbb{C}$-代数 $R$。）考虑 $\mathcal{D}_X$-代数
\[
\mathfrak{z}=\mathfrak{z}_{\mathfrak{g}}:=\mathfrak{z}_{\mathfrak{g}}(\mathcal{O})_X,
\]
它对应于交换 $(\operatorname{Aut}\mathcal{O})$-代数 $\mathfrak{z}_{\mathfrak{g}}(\mathcal{O})$（见 2.6.5、2.6.6）。
其纤维是 2.5.1 中的代数 $\mathfrak{z}_x$。
标准论证表明，当 $x\in X$ 变动时，2.5.4 中的态射 $h_x$ 定义了
$\mathcal{O}_X$-代数态射 $h: \mathfrak{z}\to\Gamma(\operatorname{Bun}_G,D')_X$。

\bdp{2.7.3}\emph{水平性定理.} $h$ 水平，即它是 $\mathcal{D}_X$-代数态射。

证明见 2.8。

\bdp{2.7.4} 置
\begin{equation}\tag{27}
\mathfrak{z}(X)=\mathfrak{z}_{\mathfrak{g}}(X):=H_\nabla(X,\mathfrak{z}).
\end{equation}
